\documentclass[10pt,a4paper]{amsart}
\usepackage[margin=19mm]{geometry}
\usepackage{amsmath,amssymb,mathtools}
\usepackage[scr=rsfs,cal=euler]{mathalfa}
\usepackage{xfrac}
\usepackage[unicode,hidelinks,bookmarks=false]{hyperref}
\newcommand{\ie}{i.e.\ }
\newcommand{\cf}{cf.\ }
\newcommand{\resp}{respectively\ }
\newcommand{\Subsection}{\subsection}
\providecommand{\nobreakdash}{\nobreak\hbox{-}}
\setlength{\emergencystretch}{2em}
\multlinegap0pt
\usepackage[shortlabels]{enumitem}
\usepackage{amssymb}
\usepackage{mathtools}
\newtheorem{proposition}{Proposition}
\newtheorem{theorem}{Theorem}
\newcommand{\Q}{\mathbb{Q}}
\title{Legendre compressions and an integrality conjecture for the H\"ormander--Bernhardsson extremal function}
\author{Khai-Hoan Nguyen-Dang}
\date{Source: arXiv:2603.24504v1 (CC BY 4.0)}
\begin{document}
\maketitle

\noindent\emph{Source and licence.} Legendre compressions and an integrality conjecture for the H\"ormander--Bernhardsson extremal function. Version arXiv:2603.24504v1 (CC BY 4.0), distributed by arXiv under the Creative Commons Attribution 4.0 International licence, http://creativecommons.org/licenses/by/4.0/, as recorded in the arXiv licence metadata for this exact version and shown on the abstract page https://arxiv.org/abs/2603.24504v1. Full licence notice: \texttt{licenses/arxiv-2603.24504-CC-BY-4.0.txt}.\par\smallskip

\noindent\textit{Editorial scope.} Counted unit: the article's theorem thm:truncation-identity (source0.tex lines 1141--1158). The article's complete original proof of it is delivered with it (source0.tex lines 1161--1223, one \texttt{proof} environment of 63 lines). No mathematics is written, repaired or re-derived: the statement and the proof are the article's own lines, in the article's own notation, and no retained line is substituted or re-flowed.  Retained uncounted as the source-local context the proof needs: lines 1014--1067.  Nothing was omitted from any retained span, so the package carries no omission record.  No retained line contains a citation, so the wrapper carries no bibliography and no bibliography entry.  No retained line contains a figure environment, an image-inclusion command or drawing code, so no image is delivered and the package declares no asset dependency.  Editorial formatting only, in addition: an AMS \texttt{amsart} preamble that restates the article's own declarations and macros used by the retained lines, namely the environments \texttt{proposition}, \texttt{theorem} on separate counters, together with the standard packages those lines need.  The retained environments print in this document as Proposition 1, Theorem 1 in order of appearance, whereas the article prints the same items with the numbering of its own full document.  The complete native source of the article as distributed in the arXiv e-print for this exact version is bundled unchanged as \texttt{sources/197152/source0.tex}.
\begin{proposition}\label{prop:truncants-determinants}
Define
\[
T_0(a,\lambda):=(-\lambda),
\]
and, for \(N\ge 1\),
\begin{equation}\label{eq:TN-def}
T_N(a,\lambda):=
\begin{pmatrix}
-\lambda & \dfrac{a}{3} & 0 & \cdots & 0 \\
-a & 2-\lambda & \dfrac{2a}{5} & \ddots & \vdots \\
0 & -\dfrac{2a}{3} & 6-\lambda & \ddots & 0 \\
\vdots & \ddots & \ddots & \ddots & \dfrac{Na}{2N+1} \\
0 & \cdots & 0 & -\dfrac{Na}{2N-1} & N(N+1)-\lambda
\end{pmatrix}.
\end{equation}
For \(N\ge 1\) and \(0\le n\le N+1\), let \(\Delta_n^{(N)}\) denote the determinant of the
trailing principal submatrix of \(T_N(a,\lambda)\) obtained by deleting the first \(n\) rows
and columns; in particular,
\[
\Delta_0^{(N)}=\det T_N(a,\lambda),
\qquad
\Delta_{N+1}^{(N)}:=1.
\]
Then the following hold.
\begin{enumerate}[label=\textup{(\roman*)}]
\item
One has
\[
\Delta_n^{(N)}=K_n^{(N)}
\qquad(1\le n\le N+1).
\]
\item
One has
\begin{equation}\label{eq:detT-K}
\det T_N(a,\lambda)=-\lambda K_1^{(N)}+\frac{a^2}{3}K_2^{(N)}.
\end{equation}
Equivalently,
\begin{equation}\label{eq:truncation-rational-identity}
\lambda q_1^{(N)}-\frac{a^2}{3}
=
-\frac{\det T_N(a,\lambda)}{K_2^{(N)}}
\end{equation}
as an identity in \(\Q(a,\lambda)\). In particular, on the locus \(K_2^{(N)}\neq 0\), the
equation
\begin{equation}\label{eq:truncation-eqn}
\lambda q_1^{(N)}=\frac{a^2}{3}
\end{equation}
is equivalent to
\begin{equation}\label{eq:truncation-det}
\det T_N(a,\lambda)=0.
\end{equation}
\end{enumerate}
\end{proposition}

\begin{theorem}[Exact finite-compression identity]\label{thm:truncation-identity}
For every \(N\ge 0\),
\[
\widehat u_{N+1}(a,\lambda)=C_{N+1}\det T_N(a,\lambda).
\]
Consequently, for every \(N\ge 0\),
\[
\widehat u_{N+1}(a,\lambda)=0
\]
if and only if
\[
\det T_N(a,\lambda)=0.
\]
If \(N\ge 1\) and \(K_2^{(N)}(a,\lambda)\neq 0\), then these equations are also equivalent to
\[
\lambda q_1^{(N)}=\frac{a^2}{3}.
\]
\end{theorem}

\begin{proof}
Set
\[
E_{-1}:=1,
\qquad
E_N:=\det T_N(a,\lambda)
\qquad(N\ge 0).
\]
Expanding along the last row of \(T_N(a,\lambda)\), we get
\[
E_N=(N(N+1)-\lambda)E_{N-1}
+\frac{N^2a^2}{(2N-1)(2N+1)}E_{N-2}
\qquad(N\ge 1),
\]
with
\[
E_0=-\lambda.
\]
Now define
\[
U_n:=C_nE_{n-1}
\qquad(n\ge 0).
\]
Then \(U_0=1\) and \(U_1=-2\lambda\). Moreover, for \(n\ge 1\),
\begin{align*}
U_{n+1}
&=C_{n+1}E_n \\
&=\frac{C_{n+1}}{C_n}\bigl(n(n+1)-\lambda\bigr)U_n
+\frac{C_{n+1}}{C_{n-1}}\frac{n^2a^2}{(2n-1)(2n+1)}U_{n-1} \\
&=\frac{2(2n+1)}{n+1}\bigl(n(n+1)-\lambda\bigr)U_n
+\frac{4n}{n+1}a^2U_{n-1} \\
&=\frac{4n+2}{n+1}\bigl(n(n+1)-\lambda\bigr)U_n
+\frac{4n}{n+1}a^2U_{n-1}.
\end{align*}
Hence \((U_n)_{n\ge 0}\) satisfies the same recurrence and the same initial conditions as
\((\widehat u_n(a,\lambda))_{n\ge 0}\). Therefore
\[
U_n=\widehat u_n(a,\lambda)
\qquad(n\ge 0).
\]
Therefore
\[
U_n=\widehat u_n(a,\lambda)
\qquad(n\ge 0),
\]
so
\[
\widehat u_{N+1}(a,\lambda)=C_{N+1}\det T_N(a,\lambda)
\qquad(N\ge 0).
\]
Since \(C_{N+1}\neq 0\), this immediately yields
\[
\widehat u_{N+1}(a,\lambda)=0
\quad\Longleftrightarrow\quad
\det T_N(a,\lambda)=0.
\]
If \(N\ge 1\) and \(K_2^{(N)}(a,\lambda)\neq 0\), then
Proposition~\ref{prop:truncants-determinants} shows that these are further equivalent to
\[
\lambda q_1^{(N)}=\frac{a^2}{3}.
\]
This completes the proof of Theorem~\ref{thm:truncation-identity}.
\end{proof}

\end{document}
